% --- Archon physics formalization source begin ---
% archon:physics
% archon:covers IPhO2026Problems/problem_IPhO_2026_3_C_2.lean
% archon:source-report reports/ipho_2026/problem_IPhO_2026_3_C_2.source.json
% archon:problem-id IPhO_2026_3
% archon:part-id C.2
% archon:previous-part IPhO_2026_3_B_1
% archon:previous-part-policy natural_language_prerequisite_only
% archon:previous-part IPhO_2026_3_C_1
% archon:previous-part-policy natural_language_prerequisite_only

\chapter{Physics problem IPhO\_2026\_3\_C\_2}
\label{ch:IPhO2026Problems_problem_IPhO_2026_3_C_2}

\paragraph{Problem source.}
The paramagnetic torus executes the Carnot refrigeration cycle
1 -> 2 -> 3 -> 4 -> 1 shown in Figure 3b in the H-versus-T plane.  T\_h and T\_c
are the hot- and cold-reservoir temperatures; Q\_h is the magnitude of heat
delivered to the hot reservoir and Q\_c is the magnitude absorbed from the cold
reservoir.  The equation of state is T*M*V = n*K*H and the isothermal heat
relation from part B may be reused.

Current subquestion:
Express M\_1 in terms of M\_2, M\_3, and M\_4.

\paragraph{Current subquestion.}
Express M\_1 in terms of M\_2, M\_3, and M\_4.

\paragraph{Recorded answer/context.}
M\_1 = sqrt(M\_2\textasciicircum{}2 - M\_3\textasciicircum{}2 + M\_4\textasciicircum{}2), taking the nonnegative magnitude.

\paragraph{Figure/image path.}
/root/proposal\_for\_physic/science-mango/ipho\_2026\_source/image/T3\_page-3.png

\paragraph{Reusable previous-part conclusions.}
\begin{itemize}
\item Source B.1. Question: At fixed temperature T, H changes from H\_i to H\_f. Find the heat Q transferred into the torus. Reusable conclusions: Q = -(mu\_0*n*K/(2*T))*(H\_f\textasciicircum{}2 - H\_i\textasciicircum{}2). Policy: natural\_language\_prerequisite\_only; do\_not\_import\_Lean\_output
\item Source C.1. Question: Label T\_h and T\_c on Figure 3b and identify the processes on which Q\_h and Q\_c are transferred. Reusable conclusions: States 1 and 4 lie at T\_h; states 2 and 3 lie at T\_c. Q\_c is absorbed on 2->3, and Q\_h is delivered on 4->1. Policy: natural\_language\_prerequisite\_only; do\_not\_import\_Lean\_output
\end{itemize}

\paragraph{Formalization target.}
create a compiling Lean file with sorry bodies at `IPhO2026Problems/problem\_IPhO\_2026\_3\_C\_2.lean`. The Lean declarations must preserve the physical quantities, dimensions or dimensional roles, figure labels, governing-law hypotheses, and final relation expressed by this problem.
Use Mathlib/Physlib names found through LeanExplore where available. If a physics API is missing, introduce faithful local abstractions rather than scalar placeholder aliases.

\begin{theorem}[Physics formalization target]
  \label{thm:physics:IPhO_2026_3_C_2:target}
  \lean{IPhO2026Problem3C2.magnetization_at_state_one}
  \uses{def:physics:IPhO_2026_3_C_2:aux001, def:physics:IPhO_2026_3_C_2:aux002, def:physics:IPhO_2026_3_C_2:aux003, def:physics:IPhO_2026_3_C_2:aux004, def:physics:IPhO_2026_3_C_2:aux005, def:physics:IPhO_2026_3_C_2:aux006, def:physics:IPhO_2026_3_C_2:aux007, def:physics:IPhO_2026_3_C_2:aux008, def:physics:IPhO_2026_3_C_2:aux009, def:physics:IPhO_2026_3_C_2:aux010, def:physics:IPhO_2026_3_C_2:aux011, def:physics:IPhO_2026_3_C_2:aux012, def:physics:IPhO_2026_3_C_2:aux013, def:physics:IPhO_2026_3_C_2:aux014, def:physics:IPhO_2026_3_C_2:aux015, def:physics:IPhO_2026_3_C_2:aux016, def:physics:IPhO_2026_3_C_2:aux017, def:physics:IPhO_2026_3_C_2:aux018, def:physics:IPhO_2026_3_C_2:aux019, def:physics:IPhO_2026_3_C_2:aux020, def:physics:IPhO_2026_3_C_2:aux021, def:physics:IPhO_2026_3_C_2:aux022, def:physics:IPhO_2026_3_C_2:aux023, def:physics:IPhO_2026_3_C_2:aux024, def:physics:IPhO_2026_3_C_2:aux025, def:physics:IPhO_2026_3_C_2:aux026, def:physics:IPhO_2026_3_C_2:aux027, lem:physics:IPhO_2026_3_C_2:aux028}
  IPhO 2026 problem 3 C.2: the magnetization at state 1, on the nonnegative magnitude branch, in terms of the magnetizations at states 2, 3, and 4.
\end{theorem}
\begin{proof}
  Use the typed geometry, governing-law, branch, and measurement interfaces listed in the declaration index.  Eliminate the intermediate readouts in their physical order and then evaluate the requested exact or uncertainty relation.
\end{proof}
\section*{Formal declaration index}
The following blocks record the typed carriers and bridge declarations used by the target.  Their dependency lists follow the declaration-level model in the covered Lean file.

\begin{definition}[Declaration magneticIntensityDimension]
  \label{def:physics:IPhO_2026_3_C_2:aux001}
  \lean{IPhO2026Problem3C2.magneticIntensityDimension}
  The SI dimension of magnetic-field strength and magnetization, A / m.
\end{definition}
\begin{proof}
  This is the typed definition of the stated physical carrier; its fields or defining equation provide the claimed data.
\end{proof}

\begin{definition}[Declaration volumeDimension]
  \label{def:physics:IPhO_2026_3_C_2:aux002}
  \lean{IPhO2026Problem3C2.volumeDimension}
  The SI dimension of volume, m³.
\end{definition}
\begin{proof}
  This is the typed definition of the stated physical carrier; its fields or defining equation provide the claimed data.
\end{proof}

\begin{definition}[Declaration curieConstantDimension]
  \label{def:physics:IPhO_2026_3_C_2:aux003}
  \lean{IPhO2026Problem3C2.curieConstantDimension}
  \uses{def:physics:IPhO_2026_3_C_2:aux002}
  The dimension of the material constant K in T M V = n K H. Since M and H have the same dimension, this is temperature times volume.
\end{definition}
\begin{proof}
  This is the typed definition of the stated physical carrier; its fields or defining equation provide the claimed data.
\end{proof}

\begin{definition}[Declaration vacuumPermeabilityDimension]
  \label{def:physics:IPhO_2026_3_C_2:aux004}
  \lean{IPhO2026Problem3C2.vacuumPermeabilityDimension}
  The SI dimension of vacuum permeability, kg m / C².
\end{definition}
\begin{proof}
  This is the typed definition of the stated physical carrier; its fields or defining equation provide the claimed data.
\end{proof}

\begin{definition}[Declaration energyDimension]
  \label{def:physics:IPhO_2026_3_C_2:aux005}
  \lean{IPhO2026Problem3C2.energyDimension}
  The SI dimension of heat/energy, kg m² / s².
\end{definition}
\begin{proof}
  This is the typed definition of the stated physical carrier; its fields or defining equation provide the claimed data.
\end{proof}

\begin{definition}[Declaration PhysicalQuantity]
  \label{def:physics:IPhO_2026_3_C_2:aux006}
  \lean{IPhO2026Problem3C2.PhysicalQuantity}
  A unit-covariant real physical quantity carrying dimension d.
\end{definition}
\begin{proof}
  This is the typed definition of the stated physical carrier; its fields or defining equation provide the claimed data.
\end{proof}

\begin{definition}[Declaration Temperature]
  \label{def:physics:IPhO_2026_3_C_2:aux007}
  \lean{IPhO2026Problem3C2.Temperature}
  \uses{def:physics:IPhO_2026_3_C_2:aux006}
  The declaration Temperature introduces a typed component of the physical model.
\end{definition}
\begin{proof}
  This is the typed definition of the stated physical carrier; its fields or defining equation provide the claimed data.
\end{proof}

\begin{definition}[Declaration Volume]
  \label{def:physics:IPhO_2026_3_C_2:aux008}
  \lean{IPhO2026Problem3C2.Volume}
  \uses{def:physics:IPhO_2026_3_C_2:aux002, def:physics:IPhO_2026_3_C_2:aux006}
  The declaration Volume introduces a typed component of the physical model.
\end{definition}
\begin{proof}
  This is the typed definition of the stated physical carrier; its fields or defining equation provide the claimed data.
\end{proof}

\begin{definition}[Declaration MagneticFieldMagnitude]
  \label{def:physics:IPhO_2026_3_C_2:aux009}
  \lean{IPhO2026Problem3C2.MagneticFieldMagnitude}
  \uses{def:physics:IPhO_2026_3_C_2:aux001, def:physics:IPhO_2026_3_C_2:aux006}
  The declaration MagneticFieldMagnitude introduces a typed component of the physical model.
\end{definition}
\begin{proof}
  This is the typed definition of the stated physical carrier; its fields or defining equation provide the claimed data.
\end{proof}

\begin{definition}[Declaration MagnetizationMagnitude]
  \label{def:physics:IPhO_2026_3_C_2:aux010}
  \lean{IPhO2026Problem3C2.MagnetizationMagnitude}
  \uses{def:physics:IPhO_2026_3_C_2:aux001, def:physics:IPhO_2026_3_C_2:aux006}
  The declaration MagnetizationMagnitude introduces a typed component of the physical model.
\end{definition}
\begin{proof}
  This is the typed definition of the stated physical carrier; its fields or defining equation provide the claimed data.
\end{proof}

\begin{definition}[Declaration CurieConstant]
  \label{def:physics:IPhO_2026_3_C_2:aux011}
  \lean{IPhO2026Problem3C2.CurieConstant}
  \uses{def:physics:IPhO_2026_3_C_2:aux003, def:physics:IPhO_2026_3_C_2:aux006}
  The declaration CurieConstant introduces a typed component of the physical model.
\end{definition}
\begin{proof}
  This is the typed definition of the stated physical carrier; its fields or defining equation provide the claimed data.
\end{proof}

\begin{definition}[Declaration VacuumPermeability]
  \label{def:physics:IPhO_2026_3_C_2:aux012}
  \lean{IPhO2026Problem3C2.VacuumPermeability}
  \uses{def:physics:IPhO_2026_3_C_2:aux004, def:physics:IPhO_2026_3_C_2:aux006}
  The declaration VacuumPermeability introduces a typed component of the physical model.
\end{definition}
\begin{proof}
  This is the typed definition of the stated physical carrier; its fields or defining equation provide the claimed data.
\end{proof}

\begin{definition}[Declaration HeatMagnitude]
  \label{def:physics:IPhO_2026_3_C_2:aux013}
  \lean{IPhO2026Problem3C2.HeatMagnitude}
  \uses{def:physics:IPhO_2026_3_C_2:aux005, def:physics:IPhO_2026_3_C_2:aux006}
  The declaration HeatMagnitude introduces a typed component of the physical model.
\end{definition}
\begin{proof}
  This is the typed definition of the stated physical carrier; its fields or defining equation provide the claimed data.
\end{proof}

\begin{definition}[Declaration siValue]
  \label{def:physics:IPhO_2026_3_C_2:aux014}
  \lean{IPhO2026Problem3C2.siValue}
  \uses{def:physics:IPhO_2026_3_C_2:aux006}
  The real coordinate of a dimensionful quantity in SI units.
\end{definition}
\begin{proof}
  This is the typed definition of the stated physical carrier; its fields or defining equation provide the claimed data.
\end{proof}

\begin{definition}[Declaration CycleVertex]
  \label{def:physics:IPhO_2026_3_C_2:aux015}
  \lean{IPhO2026Problem3C2.CycleVertex}
  The four labeled vertices in Figure 3b.
\end{definition}
\begin{proof}
  This is the typed definition of the stated physical carrier; its fields or defining equation provide the claimed data.
\end{proof}

\begin{definition}[Declaration CycleLeg]
  \label{def:physics:IPhO_2026_3_C_2:aux016}
  \lean{IPhO2026Problem3C2.CycleLeg}
  The oriented legs of the cycle 1 → 2 → 3 → 4 → 1.
\end{definition}
\begin{proof}
  This is the typed definition of the stated physical carrier; its fields or defining equation provide the claimed data.
\end{proof}

\begin{definition}[Declaration CycleLeg.initial]
  \label{def:physics:IPhO_2026_3_C_2:aux017}
  \lean{IPhO2026Problem3C2.CycleLeg.initial}
  \uses{def:physics:IPhO_2026_3_C_2:aux015, def:physics:IPhO_2026_3_C_2:aux016}
  Initial vertex of an oriented cycle leg.
\end{definition}
\begin{proof}
  This is the typed definition of the stated physical carrier; its fields or defining equation provide the claimed data.
\end{proof}

\begin{definition}[Declaration CycleLeg.final]
  \label{def:physics:IPhO_2026_3_C_2:aux018}
  \lean{IPhO2026Problem3C2.CycleLeg.final}
  \uses{def:physics:IPhO_2026_3_C_2:aux015, def:physics:IPhO_2026_3_C_2:aux016}
  Final vertex of an oriented cycle leg.
\end{definition}
\begin{proof}
  This is the typed definition of the stated physical carrier; its fields or defining equation provide the claimed data.
\end{proof}

\begin{definition}[Declaration CycleProcessKind]
  \label{def:physics:IPhO_2026_3_C_2:aux019}
  \lean{IPhO2026Problem3C2.CycleProcessKind}
  The thermodynamic kind of each leg, including the reservoir involved.
\end{definition}
\begin{proof}
  This is the typed definition of the stated physical carrier; its fields or defining equation provide the claimed data.
\end{proof}

\begin{definition}[Declaration CycleLeg.processKind]
  \label{def:physics:IPhO_2026_3_C_2:aux020}
  \lean{IPhO2026Problem3C2.CycleLeg.processKind}
  \uses{def:physics:IPhO_2026_3_C_2:aux016, def:physics:IPhO_2026_3_C_2:aux019}
  Figure/process readout for the four oriented legs of the Carnot cycle.
\end{definition}
\begin{proof}
  This is the typed definition of the stated physical carrier; its fields or defining equation provide the claimed data.
\end{proof}

\begin{definition}[Declaration ThermodynamicState]
  \label{def:physics:IPhO_2026_3_C_2:aux021}
  \lean{IPhO2026Problem3C2.ThermodynamicState}
  \uses{def:physics:IPhO_2026_3_C_2:aux007, def:physics:IPhO_2026_3_C_2:aux009, def:physics:IPhO_2026_3_C_2:aux010, def:physics:IPhO_2026_3_C_2:aux014}
  A state of the torus at one vertex of the cycle.
\end{definition}
\begin{proof}
  This is the typed definition of the stated physical carrier; its fields or defining equation provide the claimed data.
\end{proof}

\begin{definition}[Declaration ParamagneticTorus]
  \label{def:physics:IPhO_2026_3_C_2:aux022}
  \lean{IPhO2026Problem3C2.ParamagneticTorus}
  \uses{def:physics:IPhO_2026_3_C_2:aux008, def:physics:IPhO_2026_3_C_2:aux011, def:physics:IPhO_2026_3_C_2:aux012, def:physics:IPhO_2026_3_C_2:aux014}
  Fixed physical parameters of the paramagnetic torus.
\end{definition}
\begin{proof}
  This is the typed definition of the stated physical carrier; its fields or defining equation provide the claimed data.
\end{proof}

\begin{definition}[Declaration CarnotCycleData]
  \label{def:physics:IPhO_2026_3_C_2:aux023}
  \lean{IPhO2026Problem3C2.CarnotCycleData}
  \uses{def:physics:IPhO_2026_3_C_2:aux007, def:physics:IPhO_2026_3_C_2:aux013, def:physics:IPhO_2026_3_C_2:aux014, def:physics:IPhO_2026_3_C_2:aux015, def:physics:IPhO_2026_3_C_2:aux021, def:physics:IPhO_2026_3_C_2:aux022}
  The cycle data, including the reservoir temperatures, positive heat magnitudes, and the state attached to every labeled vertex.
\end{definition}
\begin{proof}
  This is the typed definition of the stated physical carrier; its fields or defining equation provide the claimed data.
\end{proof}

\begin{definition}[Declaration magnetizationSI]
  \label{def:physics:IPhO_2026_3_C_2:aux024}
  \lean{IPhO2026Problem3C2.magnetizationSI}
  \uses{def:physics:IPhO_2026_3_C_2:aux014, def:physics:IPhO_2026_3_C_2:aux015, def:physics:IPhO_2026_3_C_2:aux023}
  SI magnetization magnitude at a labeled cycle vertex.
\end{definition}
\begin{proof}
  This is the typed definition of the stated physical carrier; its fields or defining equation provide the claimed data.
\end{proof}

\begin{definition}[Declaration magneticFieldSI]
  \label{def:physics:IPhO_2026_3_C_2:aux025}
  \lean{IPhO2026Problem3C2.magneticFieldSI}
  \uses{def:physics:IPhO_2026_3_C_2:aux014, def:physics:IPhO_2026_3_C_2:aux015, def:physics:IPhO_2026_3_C_2:aux023}
  SI magnetic-field magnitude at a labeled cycle vertex.
\end{definition}
\begin{proof}
  This is the typed definition of the stated physical carrier; its fields or defining equation provide the claimed data.
\end{proof}

\begin{definition}[Declaration stateTemperatureSI]
  \label{def:physics:IPhO_2026_3_C_2:aux026}
  \lean{IPhO2026Problem3C2.stateTemperatureSI}
  \uses{def:physics:IPhO_2026_3_C_2:aux014, def:physics:IPhO_2026_3_C_2:aux015, def:physics:IPhO_2026_3_C_2:aux023}
  SI temperature at a labeled cycle vertex.
\end{definition}
\begin{proof}
  This is the typed definition of the stated physical carrier; its fields or defining equation provide the claimed data.
\end{proof}

\begin{definition}[Declaration SatisfiesParamagneticCarnotLaws]
  \label{def:physics:IPhO_2026_3_C_2:aux027}
  \lean{IPhO2026Problem3C2.SatisfiesParamagneticCarnotLaws}
  \uses{def:physics:IPhO_2026_3_C_2:aux014, def:physics:IPhO_2026_3_C_2:aux015, def:physics:IPhO_2026_3_C_2:aux023, def:physics:IPhO_2026_3_C_2:aux024, def:physics:IPhO_2026_3_C_2:aux025, def:physics:IPhO_2026_3_C_2:aux026}
  The physical laws used in part C.2. The heat-law equations preserve the orientation and sign convention from part B.1: heat transferred *into* the torus is positive. Thus Q\_c is used on 2 → 3, while the signed heat into the torus on 4 → 1 is -Q\_h.
\end{definition}
\begin{proof}
  This is the typed definition of the stated physical carrier; its fields or defining equation provide the claimed data.
\end{proof}

\begin{lemma}[Declaration magnetization\_square\_balance]
  \label{lem:physics:IPhO_2026_3_C_2:aux028}
  \lean{IPhO2026Problem3C2.magnetization_square_balance}
  \uses{def:physics:IPhO_2026_3_C_2:aux023, def:physics:IPhO_2026_3_C_2:aux024, def:physics:IPhO_2026_3_C_2:aux027}
  The algebraic square balance obtained by combining the equation of state, the two isothermal heat equations, and the reversible Carnot heat balance.
\end{lemma}
\begin{proof}
  Substitute the cited governing relations in order and use the stated positivity, orientation, and branch conditions to obtain the conclusion.
\end{proof}
% --- Archon physics formalization source end ---
